\documentclass[UTF8,11pt]{ctexart}
\usepackage[a4paper,margin=24mm]{geometry}
\usepackage{amsmath,amssymb,amsthm,mathtools,mathrsfs}
\usepackage[all]{xy}
\usepackage{xurl,hyperref,enumitem}
\usepackage{tikz}
\usetikzlibrary{arrows.meta,cd,decorations.markings,calc,patterns}
\hypersetup{hidelinks,breaklinks=true}
\setlength{\emergencystretch}{3em}
\allowdisplaybreaks
\newtheorem*{theorem}{Theorem}
\newtheorem*{thm}{Theorem}
\newtheorem*{lemma}{Lemma}
\newtheorem*{proposition}{Proposition}
\newtheorem*{prop}{Proposition}
\newtheorem*{corollary}{Corollary}
\newtheorem*{cor}{Corollary}
\newtheorem*{claim}{Claim}
\theoremstyle{definition}
\newtheorem*{definition}{Definition}
\newtheorem*{defn}{Definition}
\newtheorem*{remark}{Remark}
\newtheorem*{example}{Example}
\newcommand{\R}{\mathbb R}
\newcommand{\C}{\mathbb C}
\newcommand{\Z}{\mathbb Z}
\newcommand{\Q}{\mathbb Q}
\newcommand{\N}{\mathbb N}
\newcommand{\F}{\mathbb F}
\newcommand{\D}{\mathbb D}
\newcommand{\bB}{\mathbb B}
\newcommand{\bC}{\mathbb C}
\newcommand{\bR}{\mathbb R}
\newcommand{\bZ}{\mathbb Z}
\newcommand{\bN}{\mathbb N}
\newcommand{\bQ}{\mathbb Q}
\newcommand{\bD}{\mathbb D}
\newcommand{\bF}{\mathbb F}
\newcommand{\CP}{\mathbb{CP}}
\newcommand{\RP}{\mathbb{RP}}
\newcommand{\sO}{\mathscr O}
\newcommand{\sI}{\mathscr I}
\newcommand{\sF}{\mathscr F}
\newcommand{\sG}{\mathscr G}
\newcommand{\sH}{\mathscr H}
\newcommand{\sR}{\mathscr R}
\newcommand{\sM}{\mathscr M}
\newcommand{\sV}{\mathscr V}
\newcommand{\sU}{\mathscr U}
\DeclareMathOperator{\Hom}{Hom}
\DeclareMathOperator{\Ext}{Ext}
\DeclareMathOperator{\Tor}{Tor}
\DeclareMathOperator{\Spec}{Spec}
\DeclareMathOperator{\Proj}{Proj}
\DeclareMathOperator{\supp}{supp}
\DeclareMathOperator{\im}{im}
\DeclareMathOperator{\id}{id}
\DeclareMathOperator{\Tr}{Tr}
\DeclareMathOperator{\tr}{tr}
\DeclareMathOperator{\rank}{rank}
\DeclareMathOperator{\ord}{ord}
\DeclareMathOperator{\dist}{dist}
\DeclareMathOperator{\End}{End}
\DeclareMathOperator{\Aut}{Aut}
\DeclareMathOperator{\Gal}{Gal}
\DeclareMathOperator{\vol}{vol}
\DeclareMathOperator{\Cl}{Cl}
\DeclareMathOperator{\coker}{coker}
\DeclareMathOperator{\codim}{codim}
\DeclareMathOperator{\divergence}{div}
\DeclareMathOperator{\Ric}{Ric}
\DeclareMathOperator{\grad}{grad}
\DeclareMathOperator{\Hess}{Hess}
\newcommand{\abs}[1]{\left\lvert#1\right\rvert}
\newcommand{\sabs}[1]{\lvert#1\rvert}
\newcommand{\babs}[1]{\bigl\lvert#1\bigr\rvert}
\newcommand{\Babs}[1]{\Bigl\lvert#1\Bigr\rvert}
\newcommand{\bbabs}[1]{\biggl\lvert#1\biggr\rvert}
\newcommand{\BBabs}[1]{\Biggl\lvert#1\Biggr\rvert}
\newcommand{\norm}[1]{\left\lVert#1\right\rVert}
\newcommand{\snorm}[1]{\lVert#1\rVert}
\newcommand{\bnorm}[1]{\bigl\lVert#1\bigr\rVert}
\newcommand{\Bnorm}[1]{\Bigl\lVert#1\Bigr\rVert}
\newcommand{\bbnorm}[1]{\biggl\lVert#1\biggr\rVert}
\newcommand{\BBnorm}[1]{\Biggl\lVert#1\Biggr\rVert}
\newcommand{\widebar}[1]{\overline{#1}}
\newcommand{\nicefrac}[2]{\frac{#1}{#2}}
\newcommand{\linnprod}[2]{\langle#1,#2\rangle}
\newcommand{\myindex}[1]{#1}
\newcommand{\glsadd}[1]{}
\newcommand{\avoidbreak}{}
\newcommand{\sourceref}[1]{\textnormal{[原文：\nolinkurl{#1}]}}
\newcommand{\thmref}[1]{\sourceref{#1}}
\newcommand{\lemmaref}[1]{\sourceref{#1}}
\newcommand{\propref}[1]{\sourceref{#1}}
\newcommand{\corref}[1]{\sourceref{#1}}
\newcommand{\defnref}[1]{\sourceref{#1}}
\newcommand{\exampleref}[1]{\sourceref{#1}}
\newcommand{\exerciseref}[1]{\sourceref{#1}}
\newcommand{\figureref}[1]{\sourceref{#1}}
\newcommand{\sectionref}[1]{\sourceref{#1}}
\newcommand{\appendixref}[1]{\sourceref{#1}}
\newcommand{\stacksref}[2]{\href{https://stacks.math.columbia.edu/tag/#1}{#2 [#1]}}

\begin{document}
\section*{040\quad 线性系一般成员光滑的Bertini定理}
\subsection*{来源与计数目标}
The Stacks Project Authors, \emph{The Stacks Project}。Section 33.47，Situation 33.47.2 与 Lemma 33.47.3。使用实际取得的网页数学 TeX 正文，获取日期：2026-09-28；未取得网页构建的固定提交号，不编造版本。
\par\url{https://stacks.math.columbia.edu/tag/0FD6}
\par\url{https://stacks.math.columbia.edu/tag/0FD4}

\par\textbf{本文件只计一个问题：}由生成线丛并给出浸入的线性系证明一般零除子光滑；不另加本段未证明的几何不可约性。
\subsection*{作者命题与人类证明}

\paragraph{Situation 33.47.2.}
Let $k$ be a field, let $X$ be a scheme over $k$, let $\mathcal L$ be an invertible $\mathcal O_X$-module, let $V$ be a finite dimensional $k$-vector space, and let $\psi:V\to\Gamma(X,\mathcal L)$ be a $k$-linear map. Say $\dim(V)=r$ and we have a basis $v_1,\ldots,v_r$ of $V$. Then we obtain a ``universal divisor''
\[H_{univ}=Z(s_{univ})\subset\mathbf A^r\times_k X\]
as the zero scheme (Divisors, Definition 31.15.8) of the section
\[s_{univ}=\sum_{i=1,\ldots,r}x_i\psi(v_i)\in\Gamma(\mathbf A^r\times_k X,\operatorname{pr}_2^*\mathcal L).\]
For a field extension $k'/k$ the $k'$-points $v\in\mathbf A^r_k(k')$ correspond to vectors $(a_1,\ldots,a_r)$ of elements of $k'$. Thus we may on the one hand think of $v$ as the element $v=\sum_{i=1,\ldots,r}a_iv_i\in V\otimes_k k'$ and on the other hand we may assign to $v$ the section
\[\psi(v)=\sum_{i=1,\ldots,r}a_i\psi(v_i)\in\Gamma(X_{k'},\mathcal L|_{X_{k'}}).\]
With this notation it is clear that the fibre of $H_{univ}$ over $v\in V\otimes k'$ is the zero scheme of $\psi(v)$. In a formula:
\[H_v=H_{univ,v}=Z(\psi(v)).\]
We will denote this common value by $H_v$ as indicated. Finally, in this situation let $P$ be a property of vectors $v\in V\otimes_k k'$ for $k'/k$ an arbitrary field extension. We say $P$ holds for general $v\in V\otimes_k k'$ if there exists a nonempty Zariski open $U\subset\mathbf A^r_k$ such that if $v$ corresponds to a $k'$-point of $U$ for any $k'/k$ then $P(v)$ holds.

\begin{lemma}[33.47.3]
In Situation 33.47.2 assume
\begin{enumerate}
\item $X$ is smooth over $k$,
\item the image of $\psi:V\to\Gamma(X,\mathcal L)$ generates $\mathcal L$,
\item the corresponding morphism $\varphi_{\mathcal L,\psi}:X\to\mathbf P(V)$ is an immersion.
\end{enumerate}
Then for general $v\in V\otimes_k k'$ the scheme $H_v$ is smooth over $k'$.
\end{lemma}
\begin{proof}
(We observe that $X$ is separated and finite type as a locally closed subscheme of a projective space.) Let us use the notation introduced above the statement of the lemma. We consider the projections
\[
\xymatrix{\mathbf A^r_k\times_k X\ar[d]&H_{univ}\ar[l]\ar[ld]^p\ar[r]\ar[rd]_q&\mathbf A^r_k\times_k X\ar[d]\\X&&\mathbf A^r_k}
\]
Let $\Sigma\subset H_{univ}$ be the singular locus of the morphsm $q:H_{univ}\to\mathbf A^r_k$, i.e., the set of points where $q$ is not smooth. Then $\Sigma$ is closed because the smooth locus of a morphism is open by definition. Since the fibre of a smooth morphism is smooth, it suffices to prove $q(\Sigma)$ is contained in a proper closed subset of $\mathbf A^r_k$. Since $\Sigma$ (with reduced induced scheme structure) is a finite type scheme over $k$ it suffices to prove $\dim(\Sigma)<r$. This follows from Lemma 33.20.4. Since dimensions aren't changed by replacing $k$ by a bigger field (Morphisms, Lemma 29.29.3), we may and do assume $k$ is algebraically closed.

By dimension theory (Lemma 33.20.4), it suffices to prove that for $x\in X\setminus Z$ closed we have $p^{-1}(\{x\})\cap\Sigma$ has dimension $<r-\dim(X')$ where $X'$ is the unique irreducible component of $X$ containing $x$. As $X$ is smooth over $k$ and $x$ is a closed point we have $\dim(X')=\dim\mathfrak m_x/\mathfrak m_x^2$ (Morphisms, Lemma 29.35.12 and Algebra, Lemma 10.140.1). Thus we win if
\[\dim p^{-1}(x)\cap\Sigma<r-\dim\mathfrak m_x/\mathfrak m_x^2\]
for all $x\in X$ closed.

Since $V$ globally generated $\mathcal L$, for every irreducible component $X'$ of $X$ there is a nonempty Zariski open of $\mathbf A^r$ such that the fibres of $q$ over this open do not contain $X'$. (For example, if $x'\in X'$ is a closed point, then we can take the open corresponding to those vectors $v\in V$ such that $\psi(v)$ does not vanish at $x'$. This open will be the complement of a hyperplane in $\mathbf A^r_k$.) Let $U\subset\mathbf A^r$ be the (nonempty) intersection of these opens. Then the fibres of $q^{-1}(U)\to U$ are effective Cartier divisors on the fibres of $U\times_k X\to U$ (because a nonvanishing section of an invertible module on an integral scheme is a regular section). Hence the morphism $q^{-1}(U)\to U$ is flat by Divisors, Lemma 31.19.9. Thus for $x\in X$ closed and $v\in V=\mathbf A^r_k(k)$, if $(x,v)\in H_{univ}$, i.e., if $x\in H_v$ then $q$ is smooth at $(x,v)$ if and only if the fibre $H_v$ is smooth at $x$, see Morphisms, Lemma 29.35.14.

Consider the image $\psi(v)_x$ in the stalk $\mathcal L_x$ of the section corresponding to $v\in V$. We have
\[x\in H_v\ \Leftrightarrow\ \psi(v)_x\in\mathfrak m_x\mathcal L_x.\]
If this is true, then we have
\[H_v\text{ singular at }x\ \Leftrightarrow\ \psi(v)_x\in\mathfrak m_x^2\mathcal L_x.\]
Namely, $\psi(v)_x$ is not contained in $\mathfrak m_x^2\mathcal L_x$ $\Leftrightarrow$ the local equation for $H_v\subset X$ at $x$ is not contained in $\mathfrak m_x^2$ $\Leftrightarrow$ $\mathcal O_{H_v,x}$ is regular (Algebra, Lemma 10.106.3) $\Leftrightarrow$ $H_v$ is smooth at $x$ over $k$ (Algebra, Lemma 10.140.5). We conclude that the closed points of $p^{-1}(x)\cap\Sigma$ correspond to those $v\in V$ such that $\psi(v)_x\in\mathfrak m_x^2\mathcal L_x$. However, as $\varphi_{\mathcal L,\psi}$ is an immersion the map
\[V\longrightarrow\mathcal L_x/\mathfrak m_x^2\mathcal L_x\]
is surjective (small detail omitted). By the above, the closed points of the locus $p^{-1}(x)\cap\Sigma$ viewed as a subspace of $V$ is the kernel of this map and hence has dimension $r-\dim\mathfrak m_x/\mathfrak m_x^2-1$ as desired.
\end{proof}

\subsection*{整理说明（非作者正文）}
难点在于构造通用除子，把全局一般性问题转成奇异关联集合的维数问题，再把每个纤维识别为一阶射流映射的核。标准先修为有限型概形维数理论、光滑与正则局部环判据、有效 Cartier 除子的相对平坦性，以及浸入的局部函数与余切空间判据。原文在最后一个射流满射处写有 ``small detail omitted''；这属于所列浸入的局部判据，保留原文，不声称另附其证明。本题不依赖另一个 Bertini 定理。原文第二段出现未在本 Situation 定义的 $X\setminus Z$，其后明确写 ``for all $x\in X$ closed''；照录并在此提示，不静默修订。限定本题为原命题的无基点浸入情形，不借此前含基点的 33.47.1 增强题面。保留作者语言和维数论证，仅重排长公式。
\subsection*{可修改的简短标注（非作者正文）}
$c$：线性系一般截面的光滑性。

$a$：通用零除子；奇异关联集合；纤维维数；一阶射流与余切空间；正则局部环；相对平坦性。

\texttt{cross\_theory\_status = yes}。证明把“哪些参数给出奇异除子”表示成通用族中奇异集合的投影，并由射流线性映射的核计算其维数，返回非空开集上的光滑性。位置：33.47.3 全证。

全部标注均为非穷尽、可修改初稿；未标注不表示未使用。
\subsection*{许可与修改记录}
原数学正文作者：The Stacks Project Authors。按 GNU Free Documentation License 1.2 或后续版本复用；无不变章节、封面文字或封底文字。许可全文位于 \texttt{sources/stacks/GFDL-1.2.txt}。本条整理部分沿用同一许可。2026-09-28：助手摘录、独立排版并加入来源、范围说明与初稿标注；原作者不对本次整理背书。
\end{document}
